\documentclass[11pt]{article}
%%%%%%%%%%%%%%Include Packages%%%%%%%%%%%%%%%%%%%%%%%%%%
\usepackage{xcolor}
\usepackage{mathtools}
\usepackage[a4paper, total={6in, 8in}, margin=1.25in]{geometry}
\usepackage{amsmath}
\usepackage{amssymb}
\usepackage{paralist}
\usepackage{rsfso}
\usepackage{amsthm}
\usepackage{wasysym}
\usepackage[inline]{enumitem}
\usepackage{hyperref}
\usepackage{tocloft}
\usepackage{titlesec}
\usepackage{colortbl}
\usepackage[american]{circuitikz}
%%%%%%%%%%%%%%%%%%%%%%%%%%%%%%%%%%%%%%%%%%%%%%%%%%%%%%%%



\begin{document}
\Large \noindent \textbf{Physics 5C Discussion (Week 7)}\\
\normalsize
Department of Physics and Astronomy, UCLA\\
%\noindent TA: Jinyan Miao (jnmiu@ucla.edu)\\
%Office hours: Thursday 10-11 AM $\&$ 1-2 PM at PAB 2-429\\
%Tutoring center hours: Friday 2-3 PM at PAB 2-429 (Zoom room also available)\\
\noindent\rule{8cm}{0.4pt}\\
\textbf{Problem 1}: Jinyan in the lab has designed another circuit.\\


\noindent At the initial time, the parallel-plate capacitor \textbf{already fully charged}.\\

\noindent \textbf{Part (a)} The battery has internal resistance $r$. If both of the switches are closed, find the current going through the battery, and find the terminal voltage of the battery.\\

\noindent \textbf{Part (b)} What switch(es) should Jinyan close to keep the capacitor fully charged?\\

\noindent \textbf{Part (c)} While keeping the capacitor fully charged, Jinyan inserts a dielectric filling out the space between the two plates of the parallel plate capacitor. Find the change in voltage across the capacitor, and the change in energy stored in the capacitor, in terms of dielectric constant $\kappa$. \\

\noindent \textbf{Part (d)} If all switches are open at $t = 0$, what will happen to the capacitor?\\

\noindent \textbf{Part (e)} What will be the fastest way to discharge the capacitors, and why?


\newpage
\textbf{Problem 2}: 



\end{document}
